%%%%%%%%%%  scan idx 621  |  folio 608  %%%%%%%%%%
\textbf{Corollaire 3.4.} \uline{Soient} \(k\) \uline{un corps,} \(X\) \uline{un} \(k\)\uline{-schéma propre. On suppose que} \(\underline{\Pic}_{X_{\red}/k}\) \uline{possède un faisceau inversible universel. (Par exemple, ce sera le cas, d'après la théorie de descente, si toute composante connexe} \(X'_i\) \uline{de} \(X_{\red}\) \uline{possède un point rationnel, par exemple si} \(k\) \uline{est algébriquement clos, ou plus généralement si le} p.g.c.d. \uline{des degrés des} \(0\)\uline{-cycles de tout} \(X'_i\) \uline{est} \(1\)\uline{). Alors au-dessus de tout ouvert affine} \(V\) \uline{de} \(\Pic_{X/k}\)\uline{, il existe un faisceau inversible universel sur} \(X\times_kV\).

Finalement, on démontre I sous la forme plus précise suivante~:

\textbf{Proposition 3.5.} \uline{Soient} \(S\) \uline{un schéma noethérien intègre,} \(f:X\hookrightarrow Y\) \uline{une immersion nilpotente de} \(S\)\uline{-schémas propres. Alors il existe un ouvert non vide} \(V\) \uline{de} \(S\) \uline{tel que le morphisme}
\[
  f^*:\Pic_{Y/S}|V\longrightarrow\Pic_{X/S}|V
\]
\uline{soit représentable par un morphisme affine de présentation finie~; en fait,} \(N=\Ker f^*\) \uline{est un schéma en groupes affines à fibres unipotentes (au sens de} SGA 3 XVII).

En effet \(X\) est défini par un idéal nilpotent \(I\). Alors \(I\) possède une filtration par des idéaux \(I_0=I\supset I_1\supset\ldots\supset I_n=0\) tels que \(NI_i\subset I_{i-1}\), où \(N\) est le nilradical de \({O}_Y\) (par exemple \(I_k=N^kI\)). Soit \(X_i\) le sous-schéma de \(Y\) défini par \(I_i\). Par récurrence il suffit d'étudier les immersions \(X_i\hookrightarrow X_{i-1}\). On peut donc supposer que \(N.I=0\)~; c'est la condition i) de 3.1 et on finit en se ramenant au cas où les autres conditions (ii) à (v) sont vérifiées également. Grâce au théorème de la platitude générique (EGA IV 6.9.1) appliqué à \(\underline{O}_{X'}\), et aux \(R^if_*(I)\) et à \(E\), et compte tenu des relations de Künneth EGA III 6.9.9, on peut supposer, quitte à restreindre \(S\), que les conditions ii), iii) et iv) sont vérifiées.

